\documentclass[11pt,a4paper]{article}
\usepackage[utf8]{inputenc}
\usepackage[italian]{babel}
\usepackage{geometry}
\usepackage{graphicx}
\usepackage{hyperref}
\usepackage{listings}
\usepackage{xcolor}
\usepackage{amsmath}
\usepackage{enumitem}
\usepackage{booktabs}
\usepackage{longtable}
\usepackage{array}
\usepackage{tcolorbox}

\geometry{margin=2.5cm}

% Configurazione codice
\lstset{
    language=JavaScript,
    basicstyle=\ttfamily\small,
    keywordstyle=\color{blue}\bfseries,
    stringstyle=\color{red},
    commentstyle=\color{green!60!black},
    numberstyle=\tiny\color{gray},
    numbers=left,
    numbersep=5pt,
    frame=single,
    breaklines=true,
    showstringspaces=false,
    tabsize=2,
    escapeinside={(*@}{@*)}
}

% Box per evidenziare sezioni importanti
\newtcolorbox{importantbox}[1]{
    colback=red!5!white,
    colframe=red!75!black,
    title=#1,
    fonttitle=\bfseries
}

\newtcolorbox{warningbox}[1]{
    colback=orange!5!white,
    colframe=orange!75!black,
    title=#1,
    fonttitle=\bfseries
}

\newtcolorbox{tipbox}[1]{
    colback=blue!5!white,
    colframe=blue!75!black,
    title=#1,
    fonttitle=\bfseries
}

\title{\textbf{Analisi Componente: WorkLogService}\\
\large Documentazione Tecnica Esaustiva}
\author{Documentazione Sistema}
\date{\today}

\begin{document}

\maketitle
\tableofcontents
\newpage

\section{Storico Modifiche}

\begin{importantbox}{Stato Attuale - \today}
WorkLogService è il \textbf{Cuore Operativo} del sistema di tracciamento lavoro.
Se questo componente fallisce, l'app diventa inutile per l'utente.
\end{importantbox}

\section{Responsabilità}

\subsection{Cosa fa esattamente questo componente?}

\textbf{WorkLogService} è una classe che estende \texttt{BaseService} e gestisce i log delle ore lavorate (straordinari) degli utenti. È responsabile di:

\textbf{Funzioni principali}:
\begin{itemize}
    \item \textbf{CRUD WorkLog}: Creazione, lettura, aggiornamento ed eliminazione di log di lavoro
    \item \textbf{Validazione Relazioni}: Verifica che i progetti associati appartengano all'utente corrente (prevenzione IDOR)
    \item \textbf{Calcolo Durata}: Gestisce il calcolo automatico della durata in minuti (incluso lavoro oltre mezzanotte)
    \item \textbf{Query Specializzate}: Ricerca per mese, per progetto, aggregazioni per statistiche
    \item \textbf{Activity Tracking}: Registra automaticamente le attività di logging nel sistema di analytics
\end{itemize}

\subsection{Perché esiste?}

\begin{importantbox}{Problema Risolto}
Senza WorkLogService, ogni controller dovrebbe:
\begin{itemize}
    \item Implementare manualmente la validazione di ownership del progetto
    \item Gestire il calcolo della durata (con edge case mezzanotte)
    \item Scrivere query aggregate per statistiche
    \item Gestire la logica di ricerca per mese/progetto
    \item Registrare attività manualmente
\end{itemize}
\end{importantbox}

\textbf{Vantaggi dell'approccio}:
\begin{enumerate}
    \item \textbf{Sicurezza by Default}: Validazione ownership automatica su ogni creazione/update
    \item \textbf{Calcolo Automatico}: Durata calcolata nel Model (pre-save hook), garantisce coerenza
    \item \textbf{DRY}: Logica comune ereditata da BaseService
    \item \textbf{Activity Tracking}: Integrazione automatica con sistema analytics
    \item \textbf{Query Ottimizzate}: Indici composti per performance su query frequenti
\end{enumerate}

\section{Struttura del Codice}

\subsection{Architettura}

WorkLogService implementa il pattern \textbf{Template Method} ereditando da BaseService:

\begin{lstlisting}
[HTTP Request] POST /api/worklogs
     │
     ▼
[Controller] workLogService.create(req.tenantScope, data)
     │
     ▼
[WorkLogService.create()] 
     │
     ├─► [beforeCreate()] 
     │   ├─► validateProjectOwnership() // Verifica ownership
     │   └─► validateTimeRange() // Verifica orari validi
     │
     ├─► [BaseService.create()] 
     │   └─► [Model.save()] 
     │       └─► [pre-save hook] // Calcola durationMinutes
     │
     └─► [Activity Tracking] // Registra WORK_SESSION_LOGGED
\end{lstlisting}

\subsection{Model: WorkLog Schema}

\begin{lstlisting}
{
    projectId: ObjectId,      // Riferimento a Project (required)
    date: String,             // YYYY-MM-DD (required, regex validated)
    startTime: String,        // HH:mm (required, regex validated)
    endTime: String,          // HH:mm (required, regex validated)
    description: String,       // Max 500 caratteri
    durationMinutes: Number,   // Calcolato automaticamente (min: 0)
    user: ObjectId,           // Iniettato da multiTenancyPlugin
    createdAt: Date,          // Timestamp automatico
    updatedAt: Date           // Timestamp automatico
}
\end{lstlisting}

\textbf{Indici}:
\begin{itemize}
    \item \texttt{\{user: 1, date: -1\}}: Query per utente e data (ordinamento cronologico)
    \item \texttt{\{user: 1, projectId: 1, date: -1\}}: Query per utente, progetto e data
\end{itemize}

\subsection{Metodi Principali}

\subsubsection{Metodi Ereditati da BaseService}

\begin{itemize}
    \item \texttt{find(tenantScope, filters, options)}: Lista worklog con filtri
    \item \texttt{findById(tenantScope, id, options)}: Dettaglio singolo worklog
    \item \texttt{create(tenantScope, data)}: Crea worklog (override con activity tracking)
    \item \texttt{update(tenantScope, id, data)}: Aggiorna worklog
    \item \texttt{delete(tenantScope, id)}: Elimina worklog
    \item \texttt{search(tenantScope, searchTerm)}: Ricerca testuale su \texttt{description}
\end{itemize}

\subsubsection{Metodi Custom}

\begin{lstlisting}
// Trova worklog di un mese specifico
async findByMonth(tenantScope, year, month) {
    const monthStr = String(month).padStart(2, '0');
    const pattern = new RegExp(`^${year}-${monthStr}-`);
    return this.find(tenantScope, { date: pattern });
}

// Trova worklog di un progetto specifico
async findByProject(tenantScope, projectId, options = {}) {
    return this.find(tenantScope, { projectId }, options);
}
\end{lstlisting}

\textbf{Logica}:
\begin{itemize}
    \item \texttt{findByMonth()}: Usa regex pattern per filtrare date del mese (es. \texttt{/^2024-03-/})
    \item \texttt{findByProject()}: Filtra per \texttt{projectId} (già protetto da BaseService con filtro \texttt{user})
\end{itemize}

\subsubsection{Aggregazioni}

\begin{lstlisting}
// Calcola totali (ore, guadagni) per un periodo
async getTotalsByPeriod(tenantScope, startDate, endDate) {
    const userId = this._getUserId(tenantScope);
    
    const results = await WorkLog.aggregate([
        {
            $match: {
                user: userId,  // 🔒 SICUREZZA: Filtro esplicito
                date: { $gte: startDate, $lte: endDate },
            },
        },
        {
            $lookup: {
                from: 'projects',
                localField: 'projectId',
                foreignField: '_id',
                as: 'project',
            },
        },
        { $unwind: '$project' },
        {
            $group: {
                _id: null,
                totalMinutes: { $sum: '$durationMinutes' },
                totalEarnings: {
                    $sum: {
                        $multiply: [
                            { $divide: ['$durationMinutes', 60] },
                            '$project.rate',
                        ],
                    },
                },
                logCount: { $sum: 1 },
            },
        },
    ]);
    
    return results[0] || { totalMinutes: 0, totalEarnings: 0, logCount: 0 };
}
\end{lstlisting}

\textbf{Logica}:
\begin{itemize}
    \item \textbf{Sicurezza}: Filtro esplicito \texttt{\{user: userId\}} nella pipeline
    \item \textbf{Lookup}: Join con collection \texttt{projects} per ottenere \texttt{rate}
    \item \textbf{Calcolo Guadagni}: \texttt{(durationMinutes / 60) * project.rate}
    \item \textbf{Risultato}: Oggetto con \texttt{totalMinutes}, \texttt{totalEarnings}, \texttt{logCount}
\end{itemize}

\begin{lstlisting}
// Raggruppa worklog per data (timeline view)
async groupByDate(tenantScope, limit = 30) {
    const userId = this._getUserId(tenantScope);
    
    return WorkLog.aggregate([
        { $match: { user: userId } },  // 🔒 SICUREZZA: Filtro esplicito
        { $sort: { date: -1 } },
        {
            $group: {
                _id: '$date',
                logs: { $push: '$$ROOT' },
                totalMinutes: { $sum: '$durationMinutes' },
            },
        },
        { $sort: { _id: -1 } },
        { $limit: limit },
    ]);
}
\end{lstlisting}

\textbf{Logica}:
\begin{itemize}
    \item \textbf{Raggruppamento}: Raggruppa per campo \texttt{date}
    \item \textbf{Array Logs}: Ogni gruppo contiene array di worklog (\texttt{\$push: '\$\$ROOT'})
    \item \textbf{Totale Minuti}: Somma \texttt{durationMinutes} per data
    \item \textbf{Limit}: Limita a ultime N date (default 30)
\end{itemize}

\subsubsection{Lifecycle Hooks}

\begin{lstlisting}
// Validazione pre-creazione
async beforeCreate(tenantScope, data) {
    // 1. Validazione ownership progetto
    await this.validateProjectOwnership(tenantScope, data.projectId);
    
    // 2. Validazione orari
    this.validateTimeRange(data.startTime, data.endTime);
    
    return data;
}

// Override create per activity tracking
async create(tenantScope, data) {
    const created = await super.create(tenantScope, data);
    const userId = this._getUserId(tenantScope);
    
    // Registra attività
    await activityService.recordActivity(userId, 'WORK_SESSION_LOGGED', {
        entityId: created._id,
        category: 'work',
        metadata: {
            projectId: created.projectId,
            durationMinutes: created.durationMinutes,
            date: created.date,
            startTime: created.startTime,
            endTime: created.endTime,
        },
    });
    
    return created;
}

// Validazione pre-update
async beforeUpdate(tenantScope, id, data) {
    // Se cambia progetto, verifica ownership
    if (data.projectId) {
        await this.validateProjectOwnership(tenantScope, data.projectId);
    }
    
    // Se cambiano orari, valida range
    if (data.startTime || data.endTime) {
        const existing = await this.findById(tenantScope, id);
        const startTime = data.startTime || existing.startTime;
        const endTime = data.endTime || existing.endTime;
        this.validateTimeRange(startTime, endTime);
    }
    
    return data;
}
\end{lstlisting}

\textbf{Logica}:
\begin{itemize}
    \item \textbf{beforeCreate()}: Validazione ownership + orari prima della creazione
    \item \textbf{create()}: Override per aggiungere activity tracking dopo creazione
    \item \textbf{beforeUpdate()}: Validazione condizionale (solo se cambiano campi rilevanti)
\end{itemize}

\subsubsection{Validazioni Private}

\begin{lstlisting}
// Verifica ownership progetto
async validateProjectOwnership(tenantScope, projectId) {
    if (!projectId) {
        throw AppError.validation('Il progetto è obbligatorio');
    }

    const userId = this._getUserId(tenantScope);
    
    // 🔒 Query diretta con filtro esplicito
    const project = await Project.findOne({ 
        _id: projectId, 
        user: userId 
    });
    
    if (!project) {
        // ATTENZIONE: Non rivelare se esiste ma appartiene ad altri!
        throw AppError.notFound('Progetto');
    }
}

// Verifica range orari valido
validateTimeRange(startTime, endTime) {
    if (!startTime || !endTime) return;
    
    const [startH, startM] = startTime.split(':').map(Number);
    const [endH, endM] = endTime.split(':').map(Number);
    
    const startMinutes = startH * 60 + startM;
    const endMinutes = endH * 60 + endM;
    
    // Permetti lavoro oltre mezzanotte (es. 23:00 - 02:00)
    // ma non intervalli nulli (stesso orario)
    if (startMinutes === endMinutes) {
        throw AppError.validation(
            'L\'ora di fine deve essere diversa dall\'ora di inizio'
        );
    }
}
\end{lstlisting}

\textbf{Logica}:
\begin{itemize}
    \item \textbf{validateProjectOwnership()}: Query diretta con filtro \texttt{\{_id, user\}} per prevenire IDOR
    \item \textbf{validateTimeRange()}: Verifica che orari siano diversi (non controlla se endTime > startTime perché gestisce mezzanotte)
\end{itemize}

\section{Gestione Date e Timezone}

\subsection{Formato Date}

\begin{importantbox}{Scelta Architetturale}
Il campo \texttt{date} è una \textbf{stringa} in formato \texttt{YYYY-MM-DD}, NON un oggetto \texttt{Date} di JavaScript.
\end{importantbox}

\textbf{Perché stringa invece di Date?}
\begin{enumerate}
    \item \textbf{No Timezone Issues}: Le stringhe non hanno timezone, evitano problemi di conversione
    \item \textbf{Query Semplici}: Regex pattern per filtrare per mese (es. \texttt{/^2024-03-/})
    \item \textbf{UI Friendly}: Formato ISO standard, facile da usare in form HTML
    \item \textbf{Consistenza}: Stesso formato usato in tutto il sistema
\end{enumerate}

\textbf{Validazione}:
\begin{lstlisting}
date: { 
    type: String, 
    required: [true, 'La data è obbligatoria'],
    match: [/^\d{4}-\d{2}-\d{2}$/, 'Formato data non valido (usa YYYY-MM-DD)'],
}
\end{lstlisting}

\subsection{Formato Orari}

\begin{importantbox}{Scelta Architetturale}
I campi \texttt{startTime} e \texttt{endTime} sono \textbf{stringhe} in formato \texttt{HH:mm}, NON oggetti \texttt{Date}.
\end{importantbox}

\textbf{Perché stringa invece di Date?}
\begin{enumerate}
    \item \textbf{No Timezone}: Solo orari locali, nessuna conversione timezone
    \item \textbf{Semplicità}: Facile da validare e manipolare
    \item \textbf{UI Friendly}: Formato standard per input HTML \texttt{type="time"}
    \item \textbf{Calcolo Semplice}: Conversione diretta a minuti per calcolo durata
\end{enumerate}

\textbf{Validazione}:
\begin{lstlisting}
startTime: { 
    type: String, 
    required: [true, 'L\'ora di inizio è obbligatoria'],
    match: [/^\d{2}:\d{2}$/, 'Formato ora non valido (usa HH:mm)'],
}
\end{lstlisting}

\subsection{Calcolo Durata: Gestione Mezzanotte}

\begin{importantbox}{Edge Case Critico}
Il sistema DEVE gestire correttamente sessioni di lavoro che attraversano la mezzanotte (es. 23:00 - 02:00).
\end{importantbox}

\textbf{Implementazione nel Model (pre-save hook)}:
\begin{lstlisting}
workLogSchema.pre('save', function() {
    if (this.isModified('startTime') || this.isModified('endTime')) {
        const [startHour, startMin] = this.startTime.split(':').map(Number);
        const [endHour, endMin] = this.endTime.split(':').map(Number);
        
        let duration = (endHour * 60 + endMin) - (startHour * 60 + startMin);
        
        // Gestisce lavoro oltre mezzanotte
        if (duration < 0) {
            duration += 24 * 60;  // Aggiunge 24 ore (1440 minuti)
        }
        
        this.durationMinutes = duration;
    }
});
\end{lstlisting}

\textbf{Esempi di Calcolo}:
\begin{itemize}
    \item \textbf{23:00 - 01:00}: 
    \begin{itemize}
        \item \texttt{startMinutes = 23 * 60 + 0 = 1380}
        \item \texttt{endMinutes = 1 * 60 + 0 = 60}
        \item \texttt{duration = 60 - 1380 = -1320} (negativo!)
        \item \texttt{duration += 1440 = 120 minuti} ✅ (2 ore)
    \end{itemize}
    
    \item \textbf{09:00 - 17:30}:
    \begin{itemize}
        \item \texttt{startMinutes = 9 * 60 + 0 = 540}
        \item \texttt{endMinutes = 17 * 60 + 30 = 1050}
        \item \texttt{duration = 1050 - 540 = 510 minuti} ✅ (8.5 ore)
    \end{itemize}
    
    \item \textbf{00:00 - 00:00}:
    \begin{itemize}
        \item \texttt{duration = 0} (stesso orario)
        \item \textbf{Validazione}: \texttt{validateTimeRange()} blocca questo caso
    \end{itemize}
\end{itemize}

\textbf{Vantaggi dell'approccio}:
\begin{enumerate}
    \item \textbf{Calcolo nel Model}: Garantisce coerenza (sempre calcolato su save)
    \item \textbf{Automatico}: Non serve calcolare manualmente nel service
    \item \textbf{Edge Case Gestito}: Mezzanotte gestita correttamente
    \item \textbf{Performance}: Calcolo una sola volta (non ad ogni query)
\end{enumerate}

\section{Consistency e Sicurezza}

\subsection{Validazione Ownership Progetto}

\begin{importantbox}{Prevenzione IDOR}
La validazione di ownership è \textbf{CRITICA} per prevenire attacchi IDOR (Insecure Direct Object Reference).
\end{importantbox}

\textbf{Scenario di Attacco (senza validazione)}:
\begin{enumerate}
    \item Utente A conosce l'ID di un progetto di Utente B (es. da URL leak)
    \item Utente A crea WorkLog con \texttt{projectId} del progetto di Utente B
    \item WorkLog viene associato al progetto di Utente B
    \item Utente B vede worklog "fantasma" nel suo progetto
\end{enumerate}

\textbf{Protezione Implementata}:
\begin{lstlisting}
async validateProjectOwnership(tenantScope, projectId) {
    const userId = this._getUserId(tenantScope);
    
    // 🔒 Query diretta con filtro esplicito
    const project = await Project.findOne({ 
        _id: projectId, 
        user: userId  // ← Filtro esplicito per user
    });
    
    if (!project) {
        // Non rivelare se esiste ma appartiene ad altri!
        throw AppError.notFound('Progetto');
    }
}
\end{lstlisting}

\textbf{Quando viene chiamata?}
\begin{itemize}
    \item \textbf{beforeCreate()}: Sempre (progetto obbligatorio)
    \item \textbf{beforeUpdate()}: Solo se \texttt{data.projectId} è presente (cambio progetto)
\end{itemize}

\textbf{Perché query diretta invece di BaseService?}
\begin{itemize}
    \item \textbf{Performance}: Query diretta più veloce (no overhead BaseService)
    \item \textbf{Esplicito}: Chiaro che stiamo verificando ownership
    \item \textbf{Fail-Fast}: Errore immediato se progetto non esiste o non appartiene all'utente
\end{itemize}

\subsection{Validazione Orari}

\textbf{Regole}:
\begin{enumerate}
    \item \texttt{startTime} e \texttt{endTime} devono essere diversi (non stesso orario)
    \item \texttt{endTime} può essere minore di \texttt{startTime} (lavoro oltre mezzanotte)
    \item Formato deve essere \texttt{HH:mm} (validato da schema)
\end{enumerate}

\textbf{Implementazione}:
\begin{lstlisting}
validateTimeRange(startTime, endTime) {
    if (!startTime || !endTime) return;  // Skip se mancanti
    
    const [startH, startM] = startTime.split(':').map(Number);
    const [endH, endM] = endTime.split(':').map(Number);
    
    const startMinutes = startH * 60 + startM;
    const endMinutes = endH * 60 + endM;
    
    // Blocca stesso orario (durata zero)
    if (startMinutes === endMinutes) {
        throw AppError.validation(
            'L\'ora di fine deve essere diversa dall\'ora di inizio'
        );
    }
    // NOTA: Non blocca endTime < startTime (gestito da pre-save hook)
}
\end{lstlisting}

\section{Dipendenze}

\subsection{Dipendenze Dirette}

\begin{itemize}
    \item \textbf{BaseService}: Classe base per operazioni CRUD standardizzate
    \item \textbf{WorkLog Model}: Schema Mongoose per worklog
    \item \textbf{Project Model}: Per validazione ownership
    \item \textbf{AppError}: Per errori standardizzati
    \item \textbf{activityService}: Per registrazione attività analytics
\end{itemize}

\subsection{Dipendenze Indirette}

\begin{itemize}
    \item \textbf{multiTenancyPlugin}: Applicato a WorkLog schema per isolamento tenant
    \item \textbf{tenantContext Middleware}: Fornisce \texttt{req.tenantScope} alle routes
\end{itemize}

\section{Analisi Critica}

\subsection{Punti di Forza}

\begin{enumerate}
    \item \textbf{Sicurezza by Default}:
    \begin{itemize}
        \item Validazione ownership automatica su ogni creazione/update
        \item Query aggregate con filtri espliciti \texttt{\{user: userId\}}
        \item Prevenzione IDOR cross-tenant
    \end{itemize}
    
    \item \textbf{Calcolo Durata Automatico}:
    \begin{itemize}
        \item Calcolo nel pre-save hook garantisce coerenza
        \item Gestione corretta di lavoro oltre mezzanotte
        \item Non serve calcolare manualmente nel service
    \end{itemize}
    
    \item \textbf{Activity Tracking Integrato}:
    \begin{itemize}
        \item Registrazione automatica di \texttt{WORK\_SESSION\_LOGGED}
        \item Metadata completo per analytics
        \item Non blocca creazione se activity tracking fallisce (try-catch)
    \end{itemize}
    
    \item \textbf{Query Ottimizzate}:
    \begin{itemize}
        \item Indici composti per query frequenti
        \item Aggregazioni efficienti con \texttt{\$lookup} per join progetti
    \end{itemize}
    
    \item \textbf{DRY}:
    \begin{itemize}
        \item Eredita logica comune da BaseService
        \item Metodi custom solo per logica specifica worklog
    \end{itemize}
\end{enumerate}

\subsection{Possibili Bug/Rischi}

\begin{enumerate}
    \item \textbf{Date senza validazione formato}:
    \begin{itemize}
        \item \textbf{Stato attuale}: Validazione regex \texttt{/^\d{4}-\d{2}-\d{2}\$/} nello schema
        \item \textbf{Rischio}: Basso, ma regex non valida date invalide (es. \texttt{2024-13-45})
        \item \textbf{Miglioramento Futuro}: Validazione con \texttt{Date.parse()} o libreria date
    \end{itemize}
    
    \item \textbf{Orari senza validazione range}:
    \begin{itemize}
        \item \textbf{Stato attuale}: Regex \texttt{/^\d{2}:\d{2}\$/} valida formato ma non range
        \item \textbf{Rischio}: Accetta \texttt{25:99} (non valido)
        \item \textbf{Miglioramento Futuro}: Validazione custom che verifica \texttt{hour < 24 && minute < 60}
    \end{itemize}
    
    \item \textbf{Calcolo durata con orari invalidi}:
    \begin{itemize}
        \item \textbf{Stato attuale}: Se \texttt{startTime = "25:00"}, calcolo fallisce silenziosamente
        \item \textbf{Rischio}: \texttt{durationMinutes} potrebbe essere NaN o negativo inaspettato
        \item \textbf{Miglioramento Futuro}: Validazione range orari prima del calcolo
    \end{itemize}
    
    \item \textbf{ProjectId mancante in update}:
    \begin{itemize}
        \item \textbf{Stato attuale}: Se \texttt{data.projectId} non è presente in update, non valida ownership
        \item \textbf{Rischio}: Basso (projectId è required nello schema, ma potrebbe essere rimosso)
        \item \textbf{Miglioramento Futuro}: Considerare validazione anche se projectId non cambia (verifica che esista ancora)
    \end{itemize}
    
    \item \textbf{Aggregazione senza filtro user}:
    \begin{itemize}
        \item \textbf{Stato attuale}: ✅ \textbf{GESTITO} - Tutte le aggregazioni hanno filtro \texttt{\{user: userId\}}
        \item \textbf{Protezione}: Query dirette con filtro esplicito
    \end{itemize}
    
    \item \textbf{Activity tracking fallisce silenziosamente}:
    \begin{itemize}
        \item \textbf{Stato attuale}: Try-catch che logga errore ma non blocca creazione
        \item \textbf{Rischio}: Basso (activity tracking non è critico per funzionalità core)
        \item \textbf{Decisione}: Corretto (non bloccare creazione worklog se analytics fallisce)
    \end{itemize}
    
    \item \textbf{Regex injection in findByMonth}:
    \begin{itemize}
        \item \textbf{Stato attuale}: Pattern costruito con \texttt{new RegExp(\`\^\{year\}-\{monthStr\}-\`)}
        \item \textbf{Rischio}: Se \texttt{year} o \texttt{month} contengono caratteri speciali regex, pattern invalido
        \item \textbf{Miglioramento Futuro}: Validare che \texttt{year} e \texttt{month} siano numeri, o usare \texttt{escapeRegex()}
    \end{itemize}
    
    \item \textbf{Timezone non gestito}:
    \begin{itemize}
        \item \textbf{Stato attuale}: Date e orari sono stringhe senza timezone
        \item \textbf{Rischio}: Se utente viaggia o cambia timezone, dati potrebbero essere inconsistenti
        \item \textbf{Decisione}: Corretto per questo use case (orari locali sono sufficienti)
        \item \textbf{Miglioramento Futuro}: Considerare campo \texttt{timezone} opzionale se necessario
    \end{itemize}
\end{enumerate}

\subsection{Performance}

\begin{enumerate}
    \item \textbf{Query per mese con regex}:
    \begin{itemize}
        \item \textbf{Problema}: Regex pattern \texttt{/^2024-03-\$/} non può usare indice su campo \texttt{date}
        \item \textbf{Impatto}: Full collection scan su collezioni grandi
        \item \textbf{Miglioramento}: Usare range query \texttt{\{date: \{ \$gte: '2024-03-01', \$lte: '2024-03-31' \}\}} invece di regex
    \end{itemize}
    
    \item \textbf{Aggregazione getTotalsByPeriod}:
    \begin{itemize}
        \item \textbf{Problema}: \texttt{\$lookup} su collection \texttt{projects} per ogni worklog
        \item \textbf{Impatto}: Query lenta se molti worklog nel periodo
        \item \textbf{Miglioramento}: Considerare caching risultati per periodi frequenti
    \end{itemize}
    
    \item \textbf{Populate automatico projectId}:
    \begin{itemize}
        \item \textbf{Problema}: \texttt{populateFields: ['projectId']} popola sempre, anche quando non necessario
        \item \textbf{Impatto}: Query aggiuntive per ogni worklog
        \item \textbf{Miglioramento}: Populate selettivo solo quando necessario
    \end{itemize}
    
    \item \textbf{Indici mancanti}:
    \begin{itemize}
        \item \textbf{Stato attuale}: Indici su \texttt{\{user, date\}} e \texttt{\{user, projectId, date\}}
        \item \textbf{Copertura}: Buona per query frequenti
        \item \textbf{Miglioramento}: Considerare indice su \texttt{date} singolo se query per data senza user (raro)
    \end{itemize}
\end{enumerate}

\subsection{Miglioramenti Codice (Clean Code)}

\begin{tipbox}{Suggerimenti per Migliorare}
\end{tipbox}

\begin{enumerate}
    \item \textbf{Validazione Range Orari}:
    \begin{lstlisting}
// PRIMA
startTime: { 
    type: String, 
    match: [/^\d{2}:\d{2}$/, 'Formato ora non valido'],
}

// DOPO
startTime: { 
    type: String, 
    match: [/^\d{2}:\d{2}$/, 'Formato ora non valido'],
    validate: {
        validator: function(v) {
            const [hour, minute] = v.split(':').map(Number);
            return hour >= 0 && hour < 24 && minute >= 0 && minute < 60;
        },
        message: 'Ora non valida (usa formato HH:mm con valori validi)',
    },
}
    \end{lstlisting}
    
    \item \textbf{Validazione Date}:
    \begin{lstlisting}
// DOPO
date: { 
    type: String, 
    match: [/^\d{4}-\d{2}-\d{2}$/, 'Formato data non valido'],
    validate: {
        validator: function(v) {
            const date = new Date(v + 'T00:00:00');
            return date instanceof Date && !isNaN(date) && 
                   date.toISOString().startsWith(v);
        },
        message: 'Data non valida',
    },
}
    \end{lstlisting}
    
    \item \textbf{findByMonth con Range Query}:
    \begin{lstlisting}
// PRIMA (regex, non usa indice)
async findByMonth(tenantScope, year, month) {
    const monthStr = String(month).padStart(2, '0');
    const pattern = new RegExp(`^${year}-${monthStr}-`);
    return this.find(tenantScope, { date: pattern });
}

// DOPO (range query, usa indice)
async findByMonth(tenantScope, year, month) {
    const monthStr = String(month).padStart(2, '0');
    const startDate = `${year}-${monthStr}-01`;
    const endDate = `${year}-${monthStr}-31`;  // MongoDB accetta anche date invalide
    return this.find(tenantScope, { 
        date: { $gte: startDate, $lte: endDate } 
    });
}
    \end{lstlisting}
    
    \item \textbf{Escape Regex in findByMonth}:
    \begin{lstlisting}
// DOPO (se si mantiene regex)
async findByMonth(tenantScope, year, month) {
    // Validazione input
    if (!Number.isInteger(year) || !Number.isInteger(month) || 
        month < 1 || month > 12) {
        throw AppError.validation('Anno o mese non validi');
    }
    
    const monthStr = String(month).padStart(2, '0');
    const pattern = new RegExp(`^${year}-${monthStr}-`);
    return this.find(tenantScope, { date: pattern });
}
    \end{lstlisting}
    
    \item \textbf{TypeScript per Type Safety}:
    \begin{itemize}
        \item Convertire in TypeScript per type checking
        \item Definire interfacce per \texttt{WorkLogData}, \texttt{TotalsResult}, etc.
        \item Previene errori a compile-time
    \end{itemize}
    
    \item \textbf{Logging per Debug}:
    \begin{lstlisting}
// Aggiungere logging opzionale
async create(tenantScope, data) {
    if (process.env.DEBUG_WORKLOG === 'true') {
        console.log(`[WorkLogService.create] User: ${tenantScope.tenantId}, Data:`, data);
    }
    // ...
}
    \end{lstlisting}
\end{enumerate}

\section{UI/UX Impact}

\subsection{Velocità di Risposta}

\begin{itemize}
    \item \textbf{Query Ottimizzate}: Indici composti per query frequenti (per mese, per progetto)
    \item \textbf{Aggregazioni}: Calcolo totali efficiente con pipeline MongoDB
    \item \textbf{Populate Selettivo}: Populate solo quando necessario (configurabile)
\end{itemize}

\subsection{Sicurezza dei Dati}

\begin{importantbox}{Impatto Sicurezza}
WorkLogService garantisce che:
\begin{itemize}
    \item \textbf{Isolamento Completo}: Utente A non può mai vedere/modificare worklog di Utente B
    \item \textbf{Prevenzione IDOR}: Impossibile associare worklog a progetti di altri utenti
    \item \textbf{Validazione Ownership}: Verifica esplicita che progetti appartengano all'utente
\end{itemize}
\end{importantbox}

\subsection{Messaggi di Errore}

\begin{itemize}
    \item \textbf{Standardizzati}: Uso di \texttt{AppError} per messaggi consistenti
    \item \textbf{User-Friendly}: "Il progetto è obbligatorio" invece di "projectId required"
    \item \textbf{Sicurezza}: Non rivela se progetto esiste ma appartiene ad altri ("Progetto non trovato")
\end{itemize}

\subsection{Esperienza Utente}

\begin{itemize}
    \item \textbf{Calcolo Automatico}: Durata calcolata automaticamente, utente non deve calcolare manualmente
    \item \textbf{Gestione Mezzanotte}: Sistema gestisce correttamente sessioni che attraversano mezzanotte
    \item \textbf{Activity Tracking}: Registrazione automatica per analytics (utente non deve fare nulla)
\end{itemize}

\section{Esempi di Utilizzo}

\subsection{Esempio 1: Creazione WorkLog}

\begin{lstlisting}
// Controller
router.post('/worklogs', asyncHandler(async (req, res) => {
    const log = await workLogService.create(req.tenantScope, {
        projectId: req.body.projectId,      // Verifica ownership automatica
        date: '2024-03-15',                 // Formato YYYY-MM-DD
        startTime: '09:00',                 // Formato HH:mm
        endTime: '17:30',                   // Formato HH:mm
        description: 'Sviluppo feature X',
    });
    
    // durationMinutes calcolato automaticamente: 510 minuti (8.5 ore)
    // Activity WORK_SESSION_LOGGED registrato automaticamente
    
    res.status(201).json({ success: true, data: log });
}));
\end{lstlisting}

\subsection{Esempio 2: WorkLog Oltre Mezzanotte}

\begin{lstlisting}
// Controller
const log = await workLogService.create(req.tenantScope, {
    projectId: '507f1f77bcf86cd799439011',
    date: '2024-03-15',
    startTime: '23:00',  // Inizio notte
    endTime: '02:00',    // Fine notte (giorno successivo)
    description: 'Lavoro notturno',
});

// durationMinutes calcolato: 180 minuti (3 ore)
// Sistema gestisce correttamente il passaggio mezzanotte
\end{lstlisting}

\subsection{Esempio 3: Query per Mese}

\begin{lstlisting}
// Controller
const logs = await workLogService.findByMonth(
    req.tenantScope,
    2024,  // year
    3      // month
);

// Query eseguita: WorkLog.find({ 
//     user: userId, 
//     date: /^2024-03-/ 
// })
\end{lstlisting}

\subsection{Esempio 4: Statistiche Periodo}

\begin{lstlisting}
// Controller
const totals = await workLogService.getTotalsByPeriod(
    req.tenantScope,
    '2024-01-01',  // startDate
    '2024-03-31'   // endDate
);

// Risultato:
// {
//     totalMinutes: 14400,      // 240 ore
//     totalEarnings: 4800.00,   // €20/ora * 240 ore
//     logCount: 45
// }
\end{lstlisting}

\subsection{Esempio 5: Timeline View}

\begin{lstlisting}
// Controller
const timeline = await workLogService.groupByDate(
    req.tenantScope,
    30  // limit: ultime 30 date
);

// Risultato:
// [
//     {
//         _id: '2024-03-15',
//         logs: [ /* array di worklog del giorno */ ],
//         totalMinutes: 510
//     },
//     {
//         _id: '2024-03-14',
//         logs: [ /* ... */ ],
//         totalMinutes: 360
//     },
//     // ...
// ]
\end{lstlisting}

\section{Pattern e Principi}

\subsection{Template Method Pattern}

\textbf{Definizione}: Definisce lo scheletro di un algoritmo, delegando alcuni passi alle sottoclassi.

\textbf{Implementazione in WorkLogService}:
\begin{itemize}
    \item \textbf{Template (BaseService)}: \texttt{create()}, \texttt{update()}, \texttt{delete()}
    \item \textbf{Hook (Override)}: \texttt{beforeCreate()}, \texttt{beforeUpdate()}, \texttt{create()}
\end{itemize}

\textbf{Vantaggi}:
\begin{itemize}
    \item Codice comune in BaseService
    \item Estensibilità senza modificare codice base
    \item Validazioni custom nei hook
\end{itemize}

\subsection{Repository Pattern (Semplificato)}

\textbf{Definizione}: Astrae l'accesso ai dati, nascondendo i dettagli di implementazione.

\textbf{Implementazione in WorkLogService}:
\begin{itemize}
    \item Controller chiama \texttt{workLogService.create()}, non sa che usa MongoDB
    \item Se domani si cambia database, basta modificare BaseService
\end{itemize}

\subsection{Principi SOLID}

\begin{enumerate}
    \item \textbf{[S] Single Responsibility}: WorkLogService ha una sola responsabilità: gestione worklog
    \item \textbf{[O] Open/Closed}: Aperto per estensione (hook), chiuso per modifica
    \item \textbf{[L] Liskov Substitution}: Può sostituire BaseService (ereditarietà corretta)
    \item \textbf{[I] Interface Segregation}: Metodi granulari, usa solo quelli necessari
    \item \textbf{[D] Dependency Inversion}: Dipende da astrazione (BaseService), non da implementazione
\end{enumerate}

\section{Conclusioni}

WorkLogService è un componente \textbf{critico} per il funzionamento dell'app. Se questo fallisce, l'app diventa inutile per l'utente.

\textbf{Punti Chiave}:
\begin{itemize}
    \item ✅ Sicurezza by default con validazione ownership automatica
    \item ✅ Calcolo durata automatico con gestione mezzanotte
    \item ✅ Query ottimizzate con indici composti
    \item ✅ Activity tracking integrato per analytics
    \item ✅ DRY: eredita logica comune da BaseService
    \item ✅ Gestione date/orari come stringhe (no timezone issues)
\end{itemize}

\textbf{Raccomandazioni Future}:
\begin{enumerate}
    \item Validazione range orari nello schema (hour < 24, minute < 60)
    \item Validazione date con \texttt{Date.parse()} o libreria date
    \item \texttt{findByMonth()} con range query invece di regex (performance)
    \item Escape regex in \texttt{findByMonth()} per sicurezza
    \item Considerare TypeScript per type safety
    \item Validazione che projectId esista ancora in update (anche se non cambia)
\end{enumerate}

\textbf{Relazione con Altri Componenti}:
\begin{itemize}
    \item \textbf{BaseService}: Eredita operazioni CRUD standardizzate
    \item \textbf{TenantContext}: Fornisce \texttt{req.tenantScope} per isolamento tenant
    \item \textbf{Project Model}: Validazione ownership progetti
    \item \textbf{ActivityService}: Registrazione attività per analytics
\end{itemize}

\vspace{1cm}

\textit{Documento generato il \today}

\end{document}
